\section{Purpose of Step1b}

Step1b is the held-out generalization companion to the Step1a landscape
diagnostic. Step1a asks an in-sample visualization question: on a fixed batch of
synthetic KEP instances, how do MSE and perturbed Fenchel--Young (FY) training
trajectories move in the two-dimensional parameter space? Step1b asks the
standard ML-style question:

\begin{quote}
As the number of training graph instances increases, do the learned reward
models improve held-out kidney exchange decision quality on validation and test
graphs that were not used for training?
\end{quote}

The experiment is designed to make this question reproducible through a fixed
train/validation/test split, validation-based checkpoint selection, and paired
test-set comparisons between the two-stage baseline and decision-focused FY
training.

\paragraph{Language.}
This is a synthetic benchmark. The evaluation labels are synthetic reward
weights, not clinical outcomes. Therefore the main metric is called
\emph{synthetic-label decision gap} or \emph{oracle objective gap}, not ``true
regret''.

\section{Synthetic Benchmark and Model Class}

The default Step1b benchmark is the same controlled noisy-linear label regime as
Step1a:
\[
    w^{\mathrm{syn}}_e
    =
    \max\left\{0,\,
    (10u_e + 5c_e)(1+\xi_e)
    \right\},
\]
where \(u_e\) is the edge utility, \(c_e\) is the recipient cPRA feature, and
the multiplicative noise scale is \(\sigma=0.10\).

Both methods use the same two-feature linear probe:
\[
    \hat w_e(\theta)
    =
    \theta_1 u_e + \theta_2 c_e.
\]
The vector \([10,5]\) is the noise-free linear coefficient, not an exact
noise-corrupted label parameter.

\section{Data Splitting Design}

Step1b uses a fixed master split from
\[
\texttt{dataset/processed/step1\_noisy\_linear\_sigma010}.
\]
The default split sizes are shown in Table~\ref{tab:step1b-split}.

\begin{table}[h]
\centering
\caption{Step1b master split design.}
\label{tab:step1b-split}
\begin{tabular}{ll}
\toprule
Split component & Number of graphs \\
\midrule
Train pool & 1200 \\
Validation split & 400 \\
Held-out test split & 400 \\
\bottomrule
\end{tabular}
\end{table}

For each run, a training subset of size
\[
    n \in \{50, 200, 600, 1200\}
\]
is selected reproducibly from the 1200-graph train pool. With the same subset
seed, larger training sizes are selected by reshuffling the same train pool and
taking a longer prefix. This makes train-size comparisons stable.

\paragraph{Current archive status.}
The current formal archive used in this report contains completed runs for
\(n=50,200,600\). The \(n=1200\) run is planned and can be added to the same
figures/tables once completed.

\section{Training Methods}

\subsection{Two-stage reward-fitting baseline}

The two-stage method first trains \(\theta\) by MSE reward fitting:
\[
    \min_{\theta}
    \frac{1}{N}\sum_i
    \frac{1}{|E_i|}
    \sum_{e\in E_i}
    \left(\hat w^i_e(\theta)-w^{i,\mathrm{syn}}_e\right)^2.
\]
The checkpoint is selected by validation MSE loss. During evaluation, the KEP
solver is run with the learned predicted weights \(\hat w(\theta)\), and the
resulting decision is scored under the synthetic label weights.

\subsection{Decision-focused FY training}

The e2e model uses the same linear probe but trains with the perturbed FY
surrogate. For each training graph, the gradient estimate compares the average
perturbed optimizer solution with the synthetic-label oracle solution:
\[
    \nabla_{\theta} L_{\mathrm{FY}}(\theta)
    \approx
    X^\top
    \left(
        \bar y_\epsilon(\hat w(\theta)) - y^\star
    \right),
\]
where \(y^\star = S(w^{\mathrm{syn}})\) and
\[
    \bar y_\epsilon(\hat w)
    =
    \frac{1}{M}\sum_{m=1}^{M}
    S(\hat w + z^{(m)}).
\]

Step1b saves two e2e checkpoints:
\begin{enumerate}
    \item \textbf{e2e selected by validation decision gap}: this asks whether
    the FY trajectory produces a model with good held-out downstream decision
    quality.
    \item \textbf{e2e selected by validation FY loss}: this asks whether the FY
    surrogate itself can serve as a validation/model-selection loss.
\end{enumerate}

These two checkpoints are conceptually different and should both be reported.

\section{Evaluation Metrics}

For a graph \(i\), the synthetic-label decision gap is
\[
    G_i(\theta)
    =
    (w^{i,\mathrm{syn}})^\top y_i^\star
    -
    (w^{i,\mathrm{syn}})^\top S(\hat w^i(\theta)).
\]
The reported mean decision gap is \(\frac{1}{N}\sum_i G_i(\theta)\). The
normalized decision gap is
\[
    \frac{G_i(\theta)}
    {|(w^{i,\mathrm{syn}})^\top y_i^\star|+\varepsilon},
\]
averaged over graphs.

Step1b also reports the achieved/oracle objective ratio:
\[
    \frac{(w^{i,\mathrm{syn}})^\top S(\hat w^i(\theta))}
    {(w^{i,\mathrm{syn}})^\top y_i^\star}.
\]

For paired e2e-vs-2stage comparison, the per-graph improvement is
\[
    \Delta_i
    =
    G_i(\theta_{\mathrm{2stage}})
    -
    G_i(\theta_{\mathrm{e2e}}).
\]
Positive \(\Delta_i\) means that e2e has a smaller decision gap than 2stage on
the same test graph. Step1b reports the paired mean improvement and a bootstrap
95\% confidence interval.

\section{Formal Run Configuration}

The current formal run used the configuration in Table~\ref{tab:step1b-config}.

\begin{table}[h]
\centering
\caption{Formal Step1b configuration used for the current archive.}
\label{tab:step1b-config}
\begin{tabular}{ll}
\toprule
Item & Value \\
\midrule
Dataset & \texttt{dataset/processed/step1\_noisy\_linear\_sigma010} \\
Train sizes completed & \(50, 200, 600\) \\
Train size planned & \(1200\) \\
Split seed / subset seed / theta seed & 42 / 42 / 42 \\
2stage epochs & 500 \\
e2e epochs & 500 \\
FY perturbation scale \(\epsilon\) & 1.0 \\
FY perturbation samples \(M\) & 16 \\
Metric stride & 10 \\
Held-out test size & 400 \\
Unseen noisy-linear test size & 1000 \\
Realistic-synthetic stress test size & 2000 \\
\bottomrule
\end{tabular}
\end{table}

\section{Evaluation Sets}

Step1b currently evaluates models on three datasets.

\begin{table}[h]
\centering
\caption{Evaluation datasets used in the current Step1b report.}
\label{tab:evaluation-datasets}
\begin{tabular}{lll}
\toprule
Dataset & Size & Purpose \\
\midrule
Held-out noisy-linear test & 400 & Standard test split from the original master split \\
Unseen noisy-linear test & 1000 & Newly generated graph instances with same label regime \\
Realistic-synthetic stress test & 2000 & Different reward-label regime / distribution-shift stress test \\
\bottomrule
\end{tabular}
\end{table}

The realistic-synthetic stress test should not be described as a fully
graph-unseen test, because it may share graph structures with the original raw
batch. It is better described as a label-regime or distribution-shift stress
test.

\section{Result Tables}

\subsection{Held-out 400-graph noisy-linear test}

Table~\ref{tab:heldout400-results} shows the formal held-out test results from
the current archive. This is the standard test split associated with the
train/validation/test master split.

\begin{table}[h]
\centering
\caption{Held-out 400-graph noisy-linear test results. Positive paired
improvement means e2e has a smaller decision gap than 2stage on the same test
graphs.}
\label{tab:heldout400-results}
\resizebox{\textwidth}{!}{
\begin{tabular}{rllrrrrrr}
\toprule
Train size & Method & Selected by & Epoch & \(\theta_1\) & \(\theta_2\)
& Mean gap & Norm. gap & Paired improvement \\
\midrule
50 & 2stage & val MSE & 500 & 9.8707 & 5.1130 & 0.8480 & 0.00611 & -- \\
50 & e2e & val gap & 50 & 5.4807 & 2.8132 & 0.8220 & 0.00592 & 0.0260 \\
50 & e2e & val FY & 440 & 11.8304 & 5.0208 & 0.7587 & 0.00547 & 0.0894 \\
\midrule
200 & 2stage & val MSE & 500 & 9.8848 & 5.1038 & 0.8451 & 0.00610 & -- \\
200 & e2e & val gap & 150 & 8.8078 & 4.5582 & 0.8480 & 0.00611 & -0.0029 \\
200 & e2e & val FY & 500 & 12.2682 & 6.0649 & 0.7730 & 0.00563 & 0.0721 \\
\midrule
600 & 2stage & val MSE & 500 & 9.8899 & 5.0963 & 0.8223 & 0.00591 & -- \\
600 & e2e & val gap & 110 & 7.7838 & 4.0153 & 0.8223 & 0.00591 & 0.0000 \\
600 & e2e & val FY & 490 & 12.2328 & 5.8586 & 0.7507 & 0.00550 & 0.0716 \\
\bottomrule
\end{tabular}
}
\end{table}

\paragraph{Reading Table~\ref{tab:heldout400-results}.}
On the original 400-graph held-out split, the e2e checkpoint selected by
validation FY loss is better than the 2stage baseline for the completed
train sizes. The e2e checkpoint selected by validation decision gap is not
consistently better; it is slightly better at \(n=50\), worse at \(n=200\), and
tied at \(n=600\). This distinction is important: the FY-loss-selected model is
the stronger evidence that the FY surrogate can sometimes serve as a useful
model-selection objective on the original noisy-linear split.

\subsection{1000-graph unseen noisy-linear test}

Table~\ref{tab:unseen1000-results} evaluates the same trained models on a
separately generated unseen 1000-graph noisy-linear dataset.

\begin{table}[h]
\centering
\caption{1000-graph unseen noisy-linear test results.}
\label{tab:unseen1000-results}
\resizebox{\textwidth}{!}{
\begin{tabular}{rllrrrrrrl}
\toprule
Train size & Method & Selected by & Epoch & \(\theta_1\) & \(\theta_2\)
& Mean gap & Norm. gap & Paired improvement & 95\% CI \\
\midrule
50 & 2stage & val MSE & 500 & 9.8707 & 5.1130 & 0.7311 & 0.00531 & -- & -- \\
50 & e2e & val gap & 50 & 5.4807 & 2.8132 & 0.7352 & 0.00534 & -0.0041 & [-0.0114, 0.0006] \\
50 & e2e & val FY & 440 & 11.8304 & 5.0208 & 0.7591 & 0.00550 & -0.0280 & [-0.0878, 0.0319] \\
\midrule
200 & 2stage & val MSE & 500 & 9.8848 & 5.1038 & 0.7313 & 0.00532 & -- & -- \\
200 & e2e & val gap & 150 & 8.8078 & 4.5582 & 0.7313 & 0.00532 & 0.0000 & [0.0000, 0.0000] \\
200 & e2e & val FY & 500 & 12.2682 & 6.0649 & 0.7262 & 0.00528 & 0.0050 & [-0.0200, 0.0320] \\
\midrule
600 & 2stage & val MSE & 500 & 9.8899 & 5.0963 & 0.7309 & 0.00531 & -- & -- \\
600 & e2e & val gap & 110 & 7.7838 & 4.0153 & 0.7329 & 0.00533 & -0.0019 & [-0.0078, 0.0000] \\
600 & e2e & val FY & 490 & 12.2328 & 5.8586 & 0.7488 & 0.00545 & -0.0179 & [-0.0564, 0.0219] \\
\bottomrule
\end{tabular}
}
\end{table}

\paragraph{Reading Table~\ref{tab:unseen1000-results}.}
The 1000-graph unseen test does not replicate the stronger FY-loss-selected
improvements from the 400-graph held-out split. The 2stage baseline is stable
around mean gap \(0.731\), and e2e is tied or slightly worse except for a small
improvement at \(n=200\) selected by validation FY loss. The confidence interval
for that improvement still crosses zero. Thus, the unseen evaluation makes the
claim more cautious: the current two-feature FY model does not yet show robust
out-of-split improvement.

\subsection{Realistic-synthetic stress test}

Table~\ref{tab:realistic-results} evaluates the same noisy-linear-trained
models under a more complex realistic-synthetic label regime:
\[
    \mathrm{expected\ transplant\ count}
    \times
    \mathrm{QALY}
    \times
    \mathrm{priority\ multiplier}
    \times
    (1+\mathrm{deterministic\ noise}).
\]

\begin{table}[h]
\centering
\caption{Realistic-synthetic label-regime stress test.}
\label{tab:realistic-results}
\resizebox{\textwidth}{!}{
\begin{tabular}{rllrrrrrrrl}
\toprule
Train size & Method & Selected by & Epoch & \(\theta_1\) & \(\theta_2\)
& Mean gap & Norm. gap & Ratio & Paired improvement & 95\% CI \\
\midrule
50 & 2stage & val MSE & 500 & 9.8707 & 5.1130 & 12.1103 & 0.06603 & 0.93397 & -- & -- \\
50 & e2e & val gap & 50 & 5.4807 & 2.8132 & 12.1314 & 0.06613 & 0.93387 & -0.0210 & [-0.0750, 0.0239] \\
50 & e2e & val FY & 440 & 11.8304 & 5.0208 & 12.3439 & 0.06718 & 0.93282 & -0.2335 & [-0.4036, -0.0678] \\
\midrule
200 & 2stage & val MSE & 500 & 9.8848 & 5.1038 & 12.0837 & 0.06588 & 0.93412 & -- & -- \\
200 & e2e & val gap & 150 & 8.8078 & 4.5582 & 12.0879 & 0.06591 & 0.93409 & -0.0041 & [-0.0118, 0.0000] \\
200 & e2e & val FY & 500 & 12.2682 & 6.0649 & 12.1739 & 0.06631 & 0.93369 & -0.0902 & [-0.1674, -0.0198] \\
\midrule
600 & 2stage & val MSE & 500 & 9.8899 & 5.0963 & 12.1108 & 0.06606 & 0.93394 & -- & -- \\
600 & e2e & val gap & 110 & 7.7838 & 4.0153 & 12.1108 & 0.06606 & 0.93394 & 0.0000 & [0.0000, 0.0000] \\
600 & e2e & val FY & 490 & 12.2328 & 5.8586 & 12.1670 & 0.06618 & 0.93382 & -0.0562 & [-0.1618, 0.0471] \\
\bottomrule
\end{tabular}
}
\end{table}

\paragraph{Reading Table~\ref{tab:realistic-results}.}
Under the shifted realistic-synthetic reward labels, the noisy-linear-trained
2stage model is slightly better than or tied with the e2e models. This should
not be interpreted as a general failure of FY. It mainly shows that a
two-feature model trained on controlled noisy-linear labels does not transfer
cleanly to a different, more complex reward definition.

\section{Figure Roadmap}

The current Step1b plotting script produces 10 presentation-ready figures. Each
figure is saved as both PNG and PDF under \(\StepOneBPlotRoot\).

\begin{table}[h]
\centering
\caption{Step1b figure roadmap.}
\label{tab:figure-roadmap}
\resizebox{\textwidth}{!}{
\begin{tabular}{llp{0.52\textwidth}}
\toprule
Figure & File & What to explain \\
\midrule
1 & \texttt{01\_normalized\_decision\_gap\_by\_dataset} &
Main generalization plot. Compare normalized decision gap across train sizes and datasets. \\
2 & \texttt{02\_raw\_decision\_gap\_by\_dataset} &
Raw oracle objective gap. Useful for checking whether normalized trends are driven by scale. \\
3 & \texttt{03\_achieved\_oracle\_ratio\_by\_dataset} &
Achieved/oracle objective ratio. Higher is better and easier to interpret as percentage of oracle value. \\
4 & \texttt{04\_paired\_improvement\_over\_2stage} &
Most important comparison plot. Positive bars mean e2e beats 2stage on paired test graphs. \\
5 & \texttt{05\_selected\_parameter\_endpoints} &
Shows where selected checkpoints land in \((\theta_1,\theta_2)\) space. \\
6 & \texttt{06\_parameter\_trajectories} &
Shows training dynamics in parameter space for each train size. \\
7 & \texttt{07\_2stage\_mse\_loss\_curves} &
Checks whether MSE reward fitting converges and generalizes on validation MSE. \\
8 & \texttt{08\_e2e\_fy\_objective\_curves} &
Checks FY surrogate optimization and train/validation FY behavior. \\
9 & \texttt{09\_e2e\_decision\_gap\_curves} &
Shows whether e2e trajectory improves downstream decision gap on train/validation. \\
10 & \texttt{10\_e2e\_validation\_tradeoff} &
Directly compares validation FY objective vs validation decision gap across epochs. \\
\bottomrule
\end{tabular}
}
\end{table}

\section{Figure-by-Figure Presentation Notes}

\subsection{Main decision quality: normalized decision gap}

\begin{figure}[h]
\centering
\includegraphics[width=\textwidth]{\StepOneBPlotRoot/01_normalized_decision_gap_by_dataset.png}
\caption{Mean normalized synthetic-label decision gap across training sizes and
evaluation datasets. Lower is better.}
\label{fig:step1b-normalized-gap}
\end{figure}

\paragraph{How to explain this figure.}
Start with this figure because it is the most compact summary of generalization.
The left panel is the original held-out noisy-linear split; the middle panel is
the newly generated unseen noisy-linear test; the right panel is the
realistic-synthetic stress test. The key message is that FY-loss selection looks
promising on the original held-out split, but the advantage weakens or disappears
on the 1000-graph unseen test and does not transfer to the shifted
realistic-synthetic reward labels.

\subsection{Raw decision gap}

\begin{figure}[h]
\centering
\includegraphics[width=\textwidth]{\StepOneBPlotRoot/02_raw_decision_gap_by_dataset.png}
\caption{Mean raw oracle objective gap across datasets. Lower is better.}
\label{fig:step1b-raw-gap}
\end{figure}

\paragraph{How to explain this figure.}
Use this figure to confirm that the normalized-gap trends are not an artifact of
normalization. The raw gaps are much larger in the realistic-synthetic stress
test because the reward scale is different.

\subsection{Achieved/oracle objective ratio}

\begin{figure}[h]
\centering
\includegraphics[width=\textwidth]{\StepOneBPlotRoot/03_achieved_oracle_ratio_by_dataset.png}
\caption{Mean achieved/oracle objective ratio. Higher is better.}
\label{fig:step1b-ratio}
\end{figure}

\paragraph{How to explain this figure.}
This figure translates the decision gap into an ``oracle percentage'' view. It
is useful for communication: even when raw gaps differ across datasets, the ratio
shows how close the induced KEP solution is to the synthetic-label oracle
solution.

\subsection{Direct method performance}

\begin{figure}[h]
\centering
\includegraphics[width=\textwidth]{\StepOneBPlotRoot/04_paired_improvement_over_2stage.png}
\caption{Direct method comparison by mean normalized decision gap. Lower values
are better. Each training size shows three bars: 2stage selected by validation
MSE, e2e selected by validation decision gap, and e2e selected by validation FY
loss.}
\label{fig:step1b-direct-performance}
\end{figure}

\paragraph{How to explain this figure.}
This plot avoids the large visual uncertainty of paired-improvement confidence
intervals and instead shows the three selected models on the same absolute
performance metric. It should be read together with the paired statistics in the
tables: lower bars indicate better mean normalized decision quality, while the
paired tables still provide graph-by-graph significance evidence.

\subsection{Selected parameter endpoints}

\begin{figure}[h]
\centering
\includegraphics[width=0.82\textwidth]{\StepOneBPlotRoot/05_selected_parameter_endpoints.png}
\caption{Selected model parameters in the two-dimensional linear-probe space.}
\label{fig:step1b-selected-parameters}
\end{figure}

\paragraph{How to explain this figure.}
This figure connects Step1b back to Step1a. The 2stage checkpoints cluster near
the clean coefficient \([10,5]\). The e2e checkpoint selected by validation
decision gap may select smaller-scale parameters, while the e2e checkpoint
selected by validation FY loss often selects a larger-scale parameter vector.
This supports the interpretation that FY does not simply recover the reward
coefficient; it moves according to perturbed decision stability.

\subsection{Parameter trajectories}

\begin{figure}[h]
\centering
\includegraphics[width=\textwidth]{\StepOneBPlotRoot/06_parameter_trajectories.png}
\caption{Training trajectories in \((\theta_1,\theta_2)\) space for each
training size.}
\label{fig:step1b-trajectories}
\end{figure}

\paragraph{How to explain this figure.}
This figure shows the actual training paths. It helps explain why selected
checkpoints differ: validation decision gap and validation FY loss can choose
different epochs along the same FY trajectory. It also shows whether increasing
training size stabilizes the trajectory.

\subsection{Training and validation curves}

\begin{figure}[h]
\centering
\includegraphics[width=\textwidth]{\StepOneBPlotRoot/07_08_09_training_curves_combined.png}
\caption{Training and validation curves for the three tracked quantities:
2stage MSE, e2e perturbed FY objective, and e2e synthetic-label decision gap.
Rows correspond to metric families and columns correspond to training-set
sizes. The 2stage MSE curves are evaluated at every epoch, while the e2e FY and
decision-gap curves are evaluated every 10 epochs in the formal archive.}
\label{fig:step1b-training-curves-combined}
\end{figure}

\paragraph{How to explain this figure.}
This combined figure keeps the three training diagnostics together. The first
row is the 2stage reward-fitting sanity check: MSE should steadily fit the
synthetic labels and validation MSE should track training MSE reasonably well.
The second row shows whether the e2e FY surrogate objective is being optimized;
the validation-FY-selected checkpoint is chosen from this curve. The third row
is not the training objective, but it is the downstream decision metric used for
validation-gap checkpoint selection.

\subsection{Validation FY objective vs decision gap}

\begin{figure}[h]
\centering
\includegraphics[width=\textwidth]{\StepOneBPlotRoot/10_e2e_validation_tradeoff.png}
\caption{Relationship between validation FY objective and validation decision
gap along the e2e trajectory.}
\label{fig:step1b-tradeoff}
\end{figure}

\paragraph{How to explain this figure.}
This figure directly tests whether validation FY objective and validation
decision gap are aligned. If the point minimizing FY objective is also near a
low decision-gap point, then FY is useful not only for training but also for
model selection. If the two markers are far apart, the FY objective may still
provide useful gradients but may not be a reliable early-stopping criterion.

\section{Main Interpretation}

The current Step1b evidence is mixed but informative.

\begin{enumerate}
    \item \textbf{On the original 400-graph held-out test split, FY-loss
    selection is promising.}
    For \(n=200\) and \(n=600\), the e2e checkpoint selected by validation FY
    loss improves over 2stage with positive paired improvement. This is the
    strongest evidence that the FY surrogate can be useful on the controlled
    noisy-linear benchmark.

    \item \textbf{On the 1000-graph unseen noisy-linear dataset, the advantage is
    not robust.}
    The 2stage baseline remains stable, and e2e is tied or slightly worse in
    most cases. The only small positive e2e result is at \(n=200\) selected by
    validation FY loss, but its bootstrap interval crosses zero.

    \item \textbf{Under the realistic-synthetic stress test, noisy-linear-trained
    e2e does not transfer better than 2stage.}
    This is expected to some extent because the model was trained on a simple
    noisy-linear label but evaluated under a more complex reward label.

    \item \textbf{The two e2e selection rules answer different questions.}
    Validation decision gap asks whether the trajectory contains a good
    decision-quality point. Validation FY loss asks whether the surrogate itself
    selects a good point. Reporting both is necessary.

    \item \textbf{The current result suggests a next experimental step.}
    To make FY more competitive on unseen data, we may need better model
    capacity, stronger feature sets, warm-started FY fine-tuning, or tuning of
    \(\epsilon\), \(M\), and checkpoint rules.
\end{enumerate}

\section{What We Can and Cannot Claim}

\paragraph{Supported by current evidence.}
The FY training pipeline is now evaluated in a standard train/validation/test
style. On the original held-out noisy-linear split, the FY-loss-selected
checkpoint can outperform the two-stage MSE baseline for the completed
\(n=200\) and \(n=600\) runs.

\paragraph{Not yet supported.}
The current evidence does not show a robust FY advantage on newly generated
unseen noisy-linear graphs or under the realistic-synthetic label shift.
Therefore, we should not claim that FY generally outperforms MSE on unseen KEP
instances in the current two-feature probe setting.

\paragraph{Best summary sentence.}
\begin{quote}
Step1b upgrades Step1a from an in-sample trajectory diagnostic to a held-out
generalization study. The current results show that FY can select better
decision-quality models on the original noisy-linear split, but this advantage is
not yet robust on newly generated unseen graphs or under a shifted
realistic-synthetic reward label.
\end{quote}

\section{Next Steps}

The current report suggests the following next experimental steps:

\begin{enumerate}
    \item Complete the \(n=1200\) formal run and add it to all tables/figures.
    \item Add an MSE-warm-started FY fine-tuning method.
    \item Repeat the formal runs over multiple subset/theta seeds.
    \item Run an \(\epsilon\) and \(M\) ablation.
    \item Test richer feature sets beyond utility and recipient cPRA.
    \item Consider reporting both exact-zero-gap rate and nonzero-gap
    conditional performance, because many graphs produce identical decisions
    across methods.
\end{enumerate}
